% =====================================================================
%  BUỔI 14 — LỜI GIẢI BÀI TẬP VỀ NHÀ (từ giao-vien/02-loi-giai-bai-tap-ve-nha.md)
%  Sinh tự động bằng scripts/md2tex-luoigiai.py — không sửa tay.
%  Tài liệu giáo viên: dùng để chữa bài trên lớp và soạn đề kiểm tra.
% =====================================================================
\documentclass[11pt, a4paper]{article}
% =====================================================================
%  CAE SHANGHAI — GÓI LỆNH DÙNG CHUNG CHO TÀI LIỆU LUYỆN THI CSCA TOÁN
%  Biên dịch bằng XeLaTeX. Mọi file .tex trong buổi học \input file này.
% =====================================================================

% ---------- NGÔN NGỮ & PHÔNG CHỮ ----------
\usepackage{fontspec}
\usepackage{polyglossia}
\setmainlanguage{vietnamese}
\setotherlanguage{english}

% Chữ Latin (Anh + Việt): TeX Gyre Termes — bản clone Times New Roman,
% đi kèm TeX Live, hỗ trợ đầy đủ dấu tiếng Việt.
\setmainfont{TeX Gyre Termes}
\setsansfont{TeX Gyre Heros}

% Chữ Hán: ưu tiên Source Han Serif SC (思源宋体) — chuẩn học thuật Trung Quốc.
% Nếu máy chưa cài, tự động dùng Songti SC (宋体, có sẵn trên macOS).
% Cả hai đều thuộc họ 宋体 nên diện mạo tài liệu không đổi.
\IfFontExistsTF{Source Han Serif SC}{
  \newcommand{\cjkfamilyname}{Source Han Serif SC}
}{
  \IfFontExistsTF{Songti SC}{
    \newcommand{\cjkfamilyname}{Songti SC}
  }{
    \newcommand{\cjkfamilyname}{Noto Serif CJK SC}
  }
}

% xeCJK tự chuyển font khi gặp ký tự Hán, không cần bọc \cjk{} thủ công.
% Đây là điều kiện bắt buộc: font Latin (TeX Gyre Termes) không có glyph chữ Hán,
% nếu không khai báo xeCJK thì mọi chữ Hán viết trần sẽ bị nuốt khỏi bản in.
\usepackage{xeCJK}
\setCJKfamilyfont{caehan}{\cjkfamilyname}[Scale=0.95]
\setCJKmainfont{\cjkfamilyname}[Scale=0.95]
\setCJKsansfont{\cjkfamilyname}[Scale=0.95]
\setCJKmonofont{\cjkfamilyname}[Scale=0.95]

% \cjk{} ép dùng font Hán trong ngữ cảnh cần thiết (ví dụ trong tiêu đề in đậm
% hoặc trong bảng có định dạng riêng). Dùng \CJKfamily chứ không phải tên font,
% vì tên font chỉ là văn bản — sẽ in ra chữ thay vì đổi font.
\newcommand{\cjk}[1]{{\CJKfamily{caehan}#1}}

% ---------- TOÁN HỌC ----------
\usepackage{amsmath, amssymb, amsthm, mathtools}
\usepackage{unicode-math}
\setmathfont{TeX Gyre Termes Math}

% ---------- BỐ CỤC TRANG ----------
% Lề phải chừa vùng an toàn của letterhead:
%   banner trên chiếm 0–19mm  → top tối thiểu 2,2cm
%   banner dưới chiếm 281–297mm → bottom tối thiểu 2,2cm
% Lề trên 2,6cm (26mm) để tiêu đề đầu trang không đè banner (banner hết ở 19mm).
\usepackage[a4paper, top=2.6cm, bottom=2.4cm, left=2.5cm, right=2.0cm]{geometry}
\usepackage{fancyhdr}
\usepackage{titlesec}
\usepackage{enumitem}
\usepackage{graphicx}
\usepackage{booktabs, array, longtable, tabularx}
\usepackage{xcolor}
\usepackage[most]{tcolorbox}
\usepackage[colorlinks=true, linkcolor=cae_blue, citecolor=cae_blue, urlcolor=cae_blue]{hyperref}

% ---------- BẢNG MÀU THƯƠNG HIỆU ----------
\definecolor{cae_blue}{RGB}{0, 51, 102}
\definecolor{cae_gold}{RGB}{196, 149, 60}
\definecolor{box_bg}{RGB}{245, 247, 250}
\definecolor{box_line}{RGB}{0, 51, 102}

% ---------- TIÊU ĐỀ MỤC ----------
% Đệm đầu trang: banner letterhead chiếm 0–19mm, mà lề trên là 22mm.
% \section* đầu trang (sau \clearpage) sát lề nên phải chừa thêm, nếu không
% sẽ đè lên banner.
\titleformat{\section}
  {\normalfont\large\bfseries\sffamily\color{cae_blue}}{\thesection.}{0.6em}{}
\titleformat{\subsection}
  {\normalfont\normalsize\bfseries\sffamily\color{cae_blue!85!black}}{\thesubsection}{0.6em}{}
\titlespacing*{\section}{0pt}{14pt}{6pt}
\titlespacing*{\subsection}{0pt}{10pt}{4pt}

% Tiêu đề không đánh số ở đầu trang: \section* sau \clearpage nằm sát lề trên
% nên đè lên banner letterhead. Dùng \vspace* ngay TRƯỚC \section* không đủ,
% vì titlesec chèn khoảng cách riêng. Cách chắc chắn là đệm bằng \null.
\newcommand{\sectiondautrang}[1]{\section*{#1}}

% Dạng bài dùng \subsection*{} nên cần định dạng riêng, đậm và có gạch chân màu
\titleformat{\subsubsection}
  {\normalfont\normalsize\bfseries\sffamily\color{cae_blue!85!black}}{\thesubsubsection}{0.6em}{}

% ---------- LETTERHEAD (ảnh nền toàn trang) ----------
% Ảnh letterhead của CAE Shanghai: banner xanh trên/dưới kèm logo chìm giữa.
% Tỉ lệ ảnh 1414x2000 = 0.7070, khớp A4 (0.7071) nên chèn thẳng không méo.
% Đặt ở lớp nền dưới cùng để chữ không bị ảnh che.
% Dùng \AddToShipoutPictureBG (không có dấu *) để áp dụng cho MỌI trang.
\usepackage{eso-pic}
\AddToShipoutPictureBG{%
  \AtPageLowerLeft{%
    \includegraphics[width=\paperwidth,height=\paperheight]{images/letterhead.png}%
  }%
}

% ---------- ĐẦU / CHÂN TRANG ----------
% Letterhead đã có banner trên và dưới, nên bỏ dòng chạy để tránh chồng chữ.
% Số trang đặt trong vùng nội dung, sát trên banner dưới, lề phải.
\pagestyle{fancy}
\setlength{\headheight}{0pt}
\fancyhf{}
\renewcommand{\headrulewidth}{0pt}
\renewcommand{\footrulewidth}{0pt}
\fancyfoot[R]{\raisebox{0.9cm}{\small\color{cae_blue}\thepage}}

% ---------- MÔI TRƯỜNG NỘI DUNG ----------
% Trong tcolorbox, LaTeX đặt parindent = parskip = 0, nên các đoạn văn trong
% khung dính liền nhau. Đặt lại khoảng cách đoạn cho mọi khung.
\newcommand{\khoangcachdoan}{\setlength{\parskip}{0.5em}\setlength{\parindent}{0pt}}
\tcbset{before upper={\khoangcachdoan}}

\newtcolorbox{dinhnghia}[1][]{
  colback=box_bg, colframe=box_line, boxrule=0.6pt, arc=1pt,
  fonttitle=\bfseries, title={Định nghĩa #1}, breakable}
\newtcolorbox{dinhly}[1][]{
  colback=box_bg, colframe=cae_gold, boxrule=0.6pt, arc=1pt,
  fonttitle=\bfseries, title={Tính chất #1}, breakable}
\newtcolorbox{vidu}[1][]{
  colback=white, colframe=black!40, boxrule=0.4pt, arc=1pt,
  fonttitle=\bfseries, title={Ví dụ #1}, breakable}
\newtcolorbox{baitap}[1][]{
  colback=white, colframe=cae_blue!50, boxrule=0.4pt, arc=1pt,
  fonttitle=\bfseries, title={Câu #1}, breakable}

% Ô đáp án: viền xanh đậm hơn, dùng trong file đáp án
\newtcolorbox{dapan}[1][]{
  colback=box_bg, colframe=cae_blue, boxrule=0.6pt, arc=1pt,
  fonttitle=\bfseries, title={Câu #1}, breakable}

% ---------- TÊN CHỨNG MINH ----------
\renewcommand{\proofname}{\normalfont\bfseries Chứng minh}

% ---------- MÔN HỌC ----------
% Đổi được ở file .tex trước khi % =====================================================================
%  CAE SHANGHAI — GÓI LỆNH DÙNG CHUNG CHO TÀI LIỆU LUYỆN THI CSCA TOÁN
%  Biên dịch bằng XeLaTeX. Mọi file .tex trong buổi học \input file này.
% =====================================================================

% ---------- NGÔN NGỮ & PHÔNG CHỮ ----------
\usepackage{fontspec}
\usepackage{polyglossia}
\setmainlanguage{vietnamese}
\setotherlanguage{english}

% Chữ Latin (Anh + Việt): TeX Gyre Termes — bản clone Times New Roman,
% đi kèm TeX Live, hỗ trợ đầy đủ dấu tiếng Việt.
\setmainfont{TeX Gyre Termes}
\setsansfont{TeX Gyre Heros}

% Chữ Hán: ưu tiên Source Han Serif SC (思源宋体) — chuẩn học thuật Trung Quốc.
% Nếu máy chưa cài, tự động dùng Songti SC (宋体, có sẵn trên macOS).
% Cả hai đều thuộc họ 宋体 nên diện mạo tài liệu không đổi.
\IfFontExistsTF{Source Han Serif SC}{
  \newcommand{\cjkfamilyname}{Source Han Serif SC}
}{
  \IfFontExistsTF{Songti SC}{
    \newcommand{\cjkfamilyname}{Songti SC}
  }{
    \newcommand{\cjkfamilyname}{Noto Serif CJK SC}
  }
}

% xeCJK tự chuyển font khi gặp ký tự Hán, không cần bọc \cjk{} thủ công.
% Đây là điều kiện bắt buộc: font Latin (TeX Gyre Termes) không có glyph chữ Hán,
% nếu không khai báo xeCJK thì mọi chữ Hán viết trần sẽ bị nuốt khỏi bản in.
\usepackage{xeCJK}
\setCJKfamilyfont{caehan}{\cjkfamilyname}[Scale=0.95]
\setCJKmainfont{\cjkfamilyname}[Scale=0.95]
\setCJKsansfont{\cjkfamilyname}[Scale=0.95]
\setCJKmonofont{\cjkfamilyname}[Scale=0.95]

% \cjk{} ép dùng font Hán trong ngữ cảnh cần thiết (ví dụ trong tiêu đề in đậm
% hoặc trong bảng có định dạng riêng). Dùng \CJKfamily chứ không phải tên font,
% vì tên font chỉ là văn bản — sẽ in ra chữ thay vì đổi font.
\newcommand{\cjk}[1]{{\CJKfamily{caehan}#1}}

% ---------- TOÁN HỌC ----------
\usepackage{amsmath, amssymb, amsthm, mathtools}
\usepackage{unicode-math}
\setmathfont{TeX Gyre Termes Math}

% ---------- BỐ CỤC TRANG ----------
% Lề phải chừa vùng an toàn của letterhead:
%   banner trên chiếm 0–19mm  → top tối thiểu 2,2cm
%   banner dưới chiếm 281–297mm → bottom tối thiểu 2,2cm
% Lề trên 2,6cm (26mm) để tiêu đề đầu trang không đè banner (banner hết ở 19mm).
\usepackage[a4paper, top=2.6cm, bottom=2.4cm, left=2.5cm, right=2.0cm]{geometry}
\usepackage{fancyhdr}
\usepackage{titlesec}
\usepackage{enumitem}
\usepackage{graphicx}
\usepackage{booktabs, array, longtable, tabularx}
\usepackage{xcolor}
\usepackage[most]{tcolorbox}
\usepackage[colorlinks=true, linkcolor=cae_blue, citecolor=cae_blue, urlcolor=cae_blue]{hyperref}

% ---------- BẢNG MÀU THƯƠNG HIỆU ----------
\definecolor{cae_blue}{RGB}{0, 51, 102}
\definecolor{cae_gold}{RGB}{196, 149, 60}
\definecolor{box_bg}{RGB}{245, 247, 250}
\definecolor{box_line}{RGB}{0, 51, 102}

% ---------- TIÊU ĐỀ MỤC ----------
% Đệm đầu trang: banner letterhead chiếm 0–19mm, mà lề trên là 22mm.
% \section* đầu trang (sau \clearpage) sát lề nên phải chừa thêm, nếu không
% sẽ đè lên banner.
\titleformat{\section}
  {\normalfont\large\bfseries\sffamily\color{cae_blue}}{\thesection.}{0.6em}{}
\titleformat{\subsection}
  {\normalfont\normalsize\bfseries\sffamily\color{cae_blue!85!black}}{\thesubsection}{0.6em}{}
\titlespacing*{\section}{0pt}{14pt}{6pt}
\titlespacing*{\subsection}{0pt}{10pt}{4pt}

% Tiêu đề không đánh số ở đầu trang: \section* sau \clearpage nằm sát lề trên
% nên đè lên banner letterhead. Dùng \vspace* ngay TRƯỚC \section* không đủ,
% vì titlesec chèn khoảng cách riêng. Cách chắc chắn là đệm bằng \null.
\newcommand{\sectiondautrang}[1]{\section*{#1}}

% Dạng bài dùng \subsection*{} nên cần định dạng riêng, đậm và có gạch chân màu
\titleformat{\subsubsection}
  {\normalfont\normalsize\bfseries\sffamily\color{cae_blue!85!black}}{\thesubsubsection}{0.6em}{}

% ---------- LETTERHEAD (ảnh nền toàn trang) ----------
% Ảnh letterhead của CAE Shanghai: banner xanh trên/dưới kèm logo chìm giữa.
% Tỉ lệ ảnh 1414x2000 = 0.7070, khớp A4 (0.7071) nên chèn thẳng không méo.
% Đặt ở lớp nền dưới cùng để chữ không bị ảnh che.
% Dùng \AddToShipoutPictureBG (không có dấu *) để áp dụng cho MỌI trang.
\usepackage{eso-pic}
\AddToShipoutPictureBG{%
  \AtPageLowerLeft{%
    \includegraphics[width=\paperwidth,height=\paperheight]{images/letterhead.png}%
  }%
}

% ---------- ĐẦU / CHÂN TRANG ----------
% Letterhead đã có banner trên và dưới, nên bỏ dòng chạy để tránh chồng chữ.
% Số trang đặt trong vùng nội dung, sát trên banner dưới, lề phải.
\pagestyle{fancy}
\setlength{\headheight}{0pt}
\fancyhf{}
\renewcommand{\headrulewidth}{0pt}
\renewcommand{\footrulewidth}{0pt}
\fancyfoot[R]{\raisebox{0.9cm}{\small\color{cae_blue}\thepage}}

% ---------- MÔI TRƯỜNG NỘI DUNG ----------
% Trong tcolorbox, LaTeX đặt parindent = parskip = 0, nên các đoạn văn trong
% khung dính liền nhau. Đặt lại khoảng cách đoạn cho mọi khung.
\newcommand{\khoangcachdoan}{\setlength{\parskip}{0.5em}\setlength{\parindent}{0pt}}
\tcbset{before upper={\khoangcachdoan}}

\newtcolorbox{dinhnghia}[1][]{
  colback=box_bg, colframe=box_line, boxrule=0.6pt, arc=1pt,
  fonttitle=\bfseries, title={Định nghĩa #1}, breakable}
\newtcolorbox{dinhly}[1][]{
  colback=box_bg, colframe=cae_gold, boxrule=0.6pt, arc=1pt,
  fonttitle=\bfseries, title={Tính chất #1}, breakable}
\newtcolorbox{vidu}[1][]{
  colback=white, colframe=black!40, boxrule=0.4pt, arc=1pt,
  fonttitle=\bfseries, title={Ví dụ #1}, breakable}
\newtcolorbox{baitap}[1][]{
  colback=white, colframe=cae_blue!50, boxrule=0.4pt, arc=1pt,
  fonttitle=\bfseries, title={Câu #1}, breakable}

% Ô đáp án: viền xanh đậm hơn, dùng trong file đáp án
\newtcolorbox{dapan}[1][]{
  colback=box_bg, colframe=cae_blue, boxrule=0.6pt, arc=1pt,
  fonttitle=\bfseries, title={Câu #1}, breakable}

% ---------- TÊN CHỨNG MINH ----------
\renewcommand{\proofname}{\normalfont\bfseries Chứng minh}

% ---------- MÔN HỌC ----------
% Đổi được ở file .tex trước khi % =====================================================================
%  CAE SHANGHAI — GÓI LỆNH DÙNG CHUNG CHO TÀI LIỆU LUYỆN THI CSCA TOÁN
%  Biên dịch bằng XeLaTeX. Mọi file .tex trong buổi học \input file này.
% =====================================================================

% ---------- NGÔN NGỮ & PHÔNG CHỮ ----------
\usepackage{fontspec}
\usepackage{polyglossia}
\setmainlanguage{vietnamese}
\setotherlanguage{english}

% Chữ Latin (Anh + Việt): TeX Gyre Termes — bản clone Times New Roman,
% đi kèm TeX Live, hỗ trợ đầy đủ dấu tiếng Việt.
\setmainfont{TeX Gyre Termes}
\setsansfont{TeX Gyre Heros}

% Chữ Hán: ưu tiên Source Han Serif SC (思源宋体) — chuẩn học thuật Trung Quốc.
% Nếu máy chưa cài, tự động dùng Songti SC (宋体, có sẵn trên macOS).
% Cả hai đều thuộc họ 宋体 nên diện mạo tài liệu không đổi.
\IfFontExistsTF{Source Han Serif SC}{
  \newcommand{\cjkfamilyname}{Source Han Serif SC}
}{
  \IfFontExistsTF{Songti SC}{
    \newcommand{\cjkfamilyname}{Songti SC}
  }{
    \newcommand{\cjkfamilyname}{Noto Serif CJK SC}
  }
}

% xeCJK tự chuyển font khi gặp ký tự Hán, không cần bọc \cjk{} thủ công.
% Đây là điều kiện bắt buộc: font Latin (TeX Gyre Termes) không có glyph chữ Hán,
% nếu không khai báo xeCJK thì mọi chữ Hán viết trần sẽ bị nuốt khỏi bản in.
\usepackage{xeCJK}
\setCJKfamilyfont{caehan}{\cjkfamilyname}[Scale=0.95]
\setCJKmainfont{\cjkfamilyname}[Scale=0.95]
\setCJKsansfont{\cjkfamilyname}[Scale=0.95]
\setCJKmonofont{\cjkfamilyname}[Scale=0.95]

% \cjk{} ép dùng font Hán trong ngữ cảnh cần thiết (ví dụ trong tiêu đề in đậm
% hoặc trong bảng có định dạng riêng). Dùng \CJKfamily chứ không phải tên font,
% vì tên font chỉ là văn bản — sẽ in ra chữ thay vì đổi font.
\newcommand{\cjk}[1]{{\CJKfamily{caehan}#1}}

% ---------- TOÁN HỌC ----------
\usepackage{amsmath, amssymb, amsthm, mathtools}
\usepackage{unicode-math}
\setmathfont{TeX Gyre Termes Math}

% ---------- BỐ CỤC TRANG ----------
% Lề phải chừa vùng an toàn của letterhead:
%   banner trên chiếm 0–19mm  → top tối thiểu 2,2cm
%   banner dưới chiếm 281–297mm → bottom tối thiểu 2,2cm
% Lề trên 2,6cm (26mm) để tiêu đề đầu trang không đè banner (banner hết ở 19mm).
\usepackage[a4paper, top=2.6cm, bottom=2.4cm, left=2.5cm, right=2.0cm]{geometry}
\usepackage{fancyhdr}
\usepackage{titlesec}
\usepackage{enumitem}
\usepackage{graphicx}
\usepackage{booktabs, array, longtable, tabularx}
\usepackage{xcolor}
\usepackage[most]{tcolorbox}
\usepackage[colorlinks=true, linkcolor=cae_blue, citecolor=cae_blue, urlcolor=cae_blue]{hyperref}

% ---------- BẢNG MÀU THƯƠNG HIỆU ----------
\definecolor{cae_blue}{RGB}{0, 51, 102}
\definecolor{cae_gold}{RGB}{196, 149, 60}
\definecolor{box_bg}{RGB}{245, 247, 250}
\definecolor{box_line}{RGB}{0, 51, 102}

% ---------- TIÊU ĐỀ MỤC ----------
% Đệm đầu trang: banner letterhead chiếm 0–19mm, mà lề trên là 22mm.
% \section* đầu trang (sau \clearpage) sát lề nên phải chừa thêm, nếu không
% sẽ đè lên banner.
\titleformat{\section}
  {\normalfont\large\bfseries\sffamily\color{cae_blue}}{\thesection.}{0.6em}{}
\titleformat{\subsection}
  {\normalfont\normalsize\bfseries\sffamily\color{cae_blue!85!black}}{\thesubsection}{0.6em}{}
\titlespacing*{\section}{0pt}{14pt}{6pt}
\titlespacing*{\subsection}{0pt}{10pt}{4pt}

% Tiêu đề không đánh số ở đầu trang: \section* sau \clearpage nằm sát lề trên
% nên đè lên banner letterhead. Dùng \vspace* ngay TRƯỚC \section* không đủ,
% vì titlesec chèn khoảng cách riêng. Cách chắc chắn là đệm bằng \null.
\newcommand{\sectiondautrang}[1]{\section*{#1}}

% Dạng bài dùng \subsection*{} nên cần định dạng riêng, đậm và có gạch chân màu
\titleformat{\subsubsection}
  {\normalfont\normalsize\bfseries\sffamily\color{cae_blue!85!black}}{\thesubsubsection}{0.6em}{}

% ---------- LETTERHEAD (ảnh nền toàn trang) ----------
% Ảnh letterhead của CAE Shanghai: banner xanh trên/dưới kèm logo chìm giữa.
% Tỉ lệ ảnh 1414x2000 = 0.7070, khớp A4 (0.7071) nên chèn thẳng không méo.
% Đặt ở lớp nền dưới cùng để chữ không bị ảnh che.
% Dùng \AddToShipoutPictureBG (không có dấu *) để áp dụng cho MỌI trang.
\usepackage{eso-pic}
\AddToShipoutPictureBG{%
  \AtPageLowerLeft{%
    \includegraphics[width=\paperwidth,height=\paperheight]{images/letterhead.png}%
  }%
}

% ---------- ĐẦU / CHÂN TRANG ----------
% Letterhead đã có banner trên và dưới, nên bỏ dòng chạy để tránh chồng chữ.
% Số trang đặt trong vùng nội dung, sát trên banner dưới, lề phải.
\pagestyle{fancy}
\setlength{\headheight}{0pt}
\fancyhf{}
\renewcommand{\headrulewidth}{0pt}
\renewcommand{\footrulewidth}{0pt}
\fancyfoot[R]{\raisebox{0.9cm}{\small\color{cae_blue}\thepage}}

% ---------- MÔI TRƯỜNG NỘI DUNG ----------
% Trong tcolorbox, LaTeX đặt parindent = parskip = 0, nên các đoạn văn trong
% khung dính liền nhau. Đặt lại khoảng cách đoạn cho mọi khung.
\newcommand{\khoangcachdoan}{\setlength{\parskip}{0.5em}\setlength{\parindent}{0pt}}
\tcbset{before upper={\khoangcachdoan}}

\newtcolorbox{dinhnghia}[1][]{
  colback=box_bg, colframe=box_line, boxrule=0.6pt, arc=1pt,
  fonttitle=\bfseries, title={Định nghĩa #1}, breakable}
\newtcolorbox{dinhly}[1][]{
  colback=box_bg, colframe=cae_gold, boxrule=0.6pt, arc=1pt,
  fonttitle=\bfseries, title={Tính chất #1}, breakable}
\newtcolorbox{vidu}[1][]{
  colback=white, colframe=black!40, boxrule=0.4pt, arc=1pt,
  fonttitle=\bfseries, title={Ví dụ #1}, breakable}
\newtcolorbox{baitap}[1][]{
  colback=white, colframe=cae_blue!50, boxrule=0.4pt, arc=1pt,
  fonttitle=\bfseries, title={Câu #1}, breakable}

% Ô đáp án: viền xanh đậm hơn, dùng trong file đáp án
\newtcolorbox{dapan}[1][]{
  colback=box_bg, colframe=cae_blue, boxrule=0.6pt, arc=1pt,
  fonttitle=\bfseries, title={Câu #1}, breakable}

% ---------- TÊN CHỨNG MINH ----------
\renewcommand{\proofname}{\normalfont\bfseries Chứng minh}

% ---------- MÔN HỌC ----------
% Đổi được ở file .tex trước khi \input{cae-style} bằng \renewcommand{\monhoc}{...}
\providecommand{\monhoc}{Môn Toán}

% ---------- TRANG BÌA GỌN ----------
% #1 = loại tài liệu (TÀI LIỆU CHUẨN BỊ, TÀI LIỆU BUỔI HỌC, ...)
% #2 = số buổi
% #3 = tên buổi tiếng Việt
% #4 = tên buổi tiếng Hán
% #5 = module và tỉ trọng
\newcommand{\trangbia}[5]{%
  \begin{center}
    \vspace*{0.6cm}
    {\large\bfseries\color{cae_blue} CAE SHANGHAI}\\[0.15cm]
    {\footnotesize Tài liệu luyện thi CSCA · \monhoc}\\[0.5cm]
    {\color{cae_gold}\rule{0.55\linewidth}{0.8pt}}\\[0.35cm]
    {\small\bfseries\color{cae_blue!80!black} #1}\\[0.35cm]
    {\Large\bfseries\color{cae_blue} BUỔI #2 — #3}\\[0.2cm]
    {\large\color{cae_blue!85!black} \cjk{#4}}\\[0.25cm]
    \ifx\relax#5\relax\else{\small #5}\\[0.35cm]\fi
    {\color{cae_gold}\rule{0.55\linewidth}{0.8pt}}\\[0.3cm]
    {\footnotesize Biên soạn: Hoàng Dương, CAE Shanghai \quad·\quad \today}
  \end{center}
  \vspace{0.3cm}
}

% ---------- TRANG BÌA CHO ĐỀ KIỂM TRA ----------
% Khác \trangbia ở chỗ không có chữ "BUỔI N" — dùng cho đề thi, không phải bài giảng.
% #1 = loại tài liệu, #2 = tên kỳ thi (Việt), #3 = tên kỳ thi (Hán), #4 = thông số
\newcommand{\biade}[4]{%
  \begin{center}
    \vspace*{0.6cm}
    {\large\bfseries\color{cae_blue} CAE SHANGHAI}\\[0.15cm]
    {\footnotesize Tài liệu luyện thi CSCA · \monhoc}\\[0.5cm]
    {\color{cae_gold}\rule{0.55\linewidth}{0.8pt}}\\[0.35cm]
    {\small\bfseries\color{cae_blue!80!black} #1}\\[0.35cm]
    {\Large\bfseries\color{cae_blue} #2}\\[0.2cm]
    {\large\color{cae_blue!85!black} \cjk{#3}}\\[0.25cm]
    {\small #4}\\[0.35cm]
    {\color{cae_gold}\rule{0.55\linewidth}{0.8pt}}\\[0.3cm]
    {\footnotesize Biên soạn: Hoàng Dương, CAE Shanghai \quad·\quad \today}
  \end{center}
  \vspace{0.3cm}
}

% ---------- TIỆN ÍCH ----------
% Nhãn từ khóa in đậm ở đầu đoạn, dùng cho Từ vựng / Từ khóa nhận dạng / Chú ý
% Khoảng cách đoạn do khung tự đặt (xem \khoangcachdoan ở trên),
% ở đây chỉ cần tách đoạn và bỏ thụt đầu dòng.
\newcommand{\nhan}[1]{\par\noindent\textbf{#1.}\ \ignorespaces}

% Bảng gọn, không đánh số
\newcommand{\bang}[1]{%
  \begin{center}#1\end{center}}

% Đánh dấu ô trống cho học sinh điền
\newcommand{\chotrong}{\makebox[3.2cm]{\dotfill}}
 bằng \renewcommand{\monhoc}{...}
\providecommand{\monhoc}{Môn Toán}

% ---------- TRANG BÌA GỌN ----------
% #1 = loại tài liệu (TÀI LIỆU CHUẨN BỊ, TÀI LIỆU BUỔI HỌC, ...)
% #2 = số buổi
% #3 = tên buổi tiếng Việt
% #4 = tên buổi tiếng Hán
% #5 = module và tỉ trọng
\newcommand{\trangbia}[5]{%
  \begin{center}
    \vspace*{0.6cm}
    {\large\bfseries\color{cae_blue} CAE SHANGHAI}\\[0.15cm]
    {\footnotesize Tài liệu luyện thi CSCA · \monhoc}\\[0.5cm]
    {\color{cae_gold}\rule{0.55\linewidth}{0.8pt}}\\[0.35cm]
    {\small\bfseries\color{cae_blue!80!black} #1}\\[0.35cm]
    {\Large\bfseries\color{cae_blue} BUỔI #2 — #3}\\[0.2cm]
    {\large\color{cae_blue!85!black} \cjk{#4}}\\[0.25cm]
    \ifx\relax#5\relax\else{\small #5}\\[0.35cm]\fi
    {\color{cae_gold}\rule{0.55\linewidth}{0.8pt}}\\[0.3cm]
    {\footnotesize Biên soạn: Hoàng Dương, CAE Shanghai \quad·\quad \today}
  \end{center}
  \vspace{0.3cm}
}

% ---------- TRANG BÌA CHO ĐỀ KIỂM TRA ----------
% Khác \trangbia ở chỗ không có chữ "BUỔI N" — dùng cho đề thi, không phải bài giảng.
% #1 = loại tài liệu, #2 = tên kỳ thi (Việt), #3 = tên kỳ thi (Hán), #4 = thông số
\newcommand{\biade}[4]{%
  \begin{center}
    \vspace*{0.6cm}
    {\large\bfseries\color{cae_blue} CAE SHANGHAI}\\[0.15cm]
    {\footnotesize Tài liệu luyện thi CSCA · \monhoc}\\[0.5cm]
    {\color{cae_gold}\rule{0.55\linewidth}{0.8pt}}\\[0.35cm]
    {\small\bfseries\color{cae_blue!80!black} #1}\\[0.35cm]
    {\Large\bfseries\color{cae_blue} #2}\\[0.2cm]
    {\large\color{cae_blue!85!black} \cjk{#3}}\\[0.25cm]
    {\small #4}\\[0.35cm]
    {\color{cae_gold}\rule{0.55\linewidth}{0.8pt}}\\[0.3cm]
    {\footnotesize Biên soạn: Hoàng Dương, CAE Shanghai \quad·\quad \today}
  \end{center}
  \vspace{0.3cm}
}

% ---------- TIỆN ÍCH ----------
% Nhãn từ khóa in đậm ở đầu đoạn, dùng cho Từ vựng / Từ khóa nhận dạng / Chú ý
% Khoảng cách đoạn do khung tự đặt (xem \khoangcachdoan ở trên),
% ở đây chỉ cần tách đoạn và bỏ thụt đầu dòng.
\newcommand{\nhan}[1]{\par\noindent\textbf{#1.}\ \ignorespaces}

% Bảng gọn, không đánh số
\newcommand{\bang}[1]{%
  \begin{center}#1\end{center}}

% Đánh dấu ô trống cho học sinh điền
\newcommand{\chotrong}{\makebox[3.2cm]{\dotfill}}
 bằng \renewcommand{\monhoc}{...}
\providecommand{\monhoc}{Môn Toán}

% ---------- TRANG BÌA GỌN ----------
% #1 = loại tài liệu (TÀI LIỆU CHUẨN BỊ, TÀI LIỆU BUỔI HỌC, ...)
% #2 = số buổi
% #3 = tên buổi tiếng Việt
% #4 = tên buổi tiếng Hán
% #5 = module và tỉ trọng
\newcommand{\trangbia}[5]{%
  \begin{center}
    \vspace*{0.6cm}
    {\large\bfseries\color{cae_blue} CAE SHANGHAI}\\[0.15cm]
    {\footnotesize Tài liệu luyện thi CSCA · \monhoc}\\[0.5cm]
    {\color{cae_gold}\rule{0.55\linewidth}{0.8pt}}\\[0.35cm]
    {\small\bfseries\color{cae_blue!80!black} #1}\\[0.35cm]
    {\Large\bfseries\color{cae_blue} BUỔI #2 — #3}\\[0.2cm]
    {\large\color{cae_blue!85!black} \cjk{#4}}\\[0.25cm]
    \ifx\relax#5\relax\else{\small #5}\\[0.35cm]\fi
    {\color{cae_gold}\rule{0.55\linewidth}{0.8pt}}\\[0.3cm]
    {\footnotesize Biên soạn: Hoàng Dương, CAE Shanghai \quad·\quad \today}
  \end{center}
  \vspace{0.3cm}
}

% ---------- TRANG BÌA CHO ĐỀ KIỂM TRA ----------
% Khác \trangbia ở chỗ không có chữ "BUỔI N" — dùng cho đề thi, không phải bài giảng.
% #1 = loại tài liệu, #2 = tên kỳ thi (Việt), #3 = tên kỳ thi (Hán), #4 = thông số
\newcommand{\biade}[4]{%
  \begin{center}
    \vspace*{0.6cm}
    {\large\bfseries\color{cae_blue} CAE SHANGHAI}\\[0.15cm]
    {\footnotesize Tài liệu luyện thi CSCA · \monhoc}\\[0.5cm]
    {\color{cae_gold}\rule{0.55\linewidth}{0.8pt}}\\[0.35cm]
    {\small\bfseries\color{cae_blue!80!black} #1}\\[0.35cm]
    {\Large\bfseries\color{cae_blue} #2}\\[0.2cm]
    {\large\color{cae_blue!85!black} \cjk{#3}}\\[0.25cm]
    {\small #4}\\[0.35cm]
    {\color{cae_gold}\rule{0.55\linewidth}{0.8pt}}\\[0.3cm]
    {\footnotesize Biên soạn: Hoàng Dương, CAE Shanghai \quad·\quad \today}
  \end{center}
  \vspace{0.3cm}
}

% ---------- TIỆN ÍCH ----------
% Nhãn từ khóa in đậm ở đầu đoạn, dùng cho Từ vựng / Từ khóa nhận dạng / Chú ý
% Khoảng cách đoạn do khung tự đặt (xem \khoangcachdoan ở trên),
% ở đây chỉ cần tách đoạn và bỏ thụt đầu dòng.
\newcommand{\nhan}[1]{\par\noindent\textbf{#1.}\ \ignorespaces}

% Bảng gọn, không đánh số
\newcommand{\bang}[1]{%
  \begin{center}#1\end{center}}

% Đánh dấu ô trống cho học sinh điền
\newcommand{\chotrong}{\makebox[3.2cm]{\dotfill}}


\begin{document}

\trangbia{LỜI GIẢI BÀI TẬP VỀ NHÀ}{14}{Xác suất}{概率}{Tài liệu giáo viên \quad·\quad 40 câu}

\subsection*{Nhóm A — Xác suất cổ điển và đếm bằng tổ hợp}

\begin{dapan}[1 — Đáp án D]
\textbf{Giải.} Tổng số cách chọn $3$ học sinh từ $8$ học sinh:

$$C\_8^3 = \frac{8 \cdot 7 \cdot 6}{6} = 56$$

Số cách thuận lợi: chọn $1$ nam từ $3$ nam và $2$ nữ từ $5$ nữ.

$$C\_3^1 \cdot C\_5^2 = 3 \cdot 10 = 30$$

$$P = \frac{30}{56} = \frac{15}{28}$$

\textbf{Bẫy.} Đáp án A là $\dfrac{15}{56}$, tức quên nhân với $C_3^1$ ở tử số. Đáp án C là $\dfrac{13}{56}$, không ứng với cách đếm nào.
\end{dapan}

\begin{dapan}[2 — Đáp án D]
\textbf{Giải.} Lập bảng số liệu theo hai chiều giới tính và thị lực:

\bang{
\begin{tabularx}{\linewidth}{@{\hspace{6pt}}Xccc@{\hspace{6pt}}}
\toprule
\textbf{Nhóm} & \textbf{Nam} & \textbf{Nữ} & \textbf{Tổng} \\
\midrule
Thị lực bình thường & 150 & 120 & 270 \\
Thị lực bất thường & 50 & 30 & 80 \\
Tổng & 200 & 150 & 350 \\
\bottomrule
\end{tabularx}}

Tổng số học sinh là $350$. Số nam có thị lực bình thường là $150$.

$$P = \frac{150}{350} = \frac{3}{7}$$

\textbf{Bẫy.} Đáp án B là $\dfrac{150}{200}$, tức lấy mẫu số là số nam chứ không phải cả trường. Đáp án A là $\dfrac{50}{350}$, tức đếm nhầm ô nam bị bất thường.
\end{dapan}

\begin{dapan}[3 — Đáp án B]
\textbf{Giải.} Trong $10$ số từ $1$ đến $10$ có $5$ số lẻ $\{1,3,5,7,9\}$ và $5$ số chẵn $\{2,4,6,8,10\}$. Tổng hai số là số chẵn khi hai số cùng tính chẵn lẻ.

Tổng số cách rút $2$ thẻ:

$$C\_{10}^2 = 45$$

Số cách thuận lợi:

$$C\_5^2 + C\_5^2 = 10 + 10 = 20$$

$$P = \frac{20}{45} = \frac{4}{9}$$

\textbf{Cách nhanh.} Tử số luôn là $2 C_5^2$ khi hai nhóm có số phần tử bằng nhau. Không cần tính riêng từng nhóm.
\end{dapan}

\begin{dapan}[4 — Đáp án C]
\textbf{Giải.} Hai con xúc xắc cho $6 \times 6 = 36$ kết quả. Các cặp có tổng bằng $8$:

$$(2,6), (3,5), (4,4), (5,3), (6,2)$$

Có $5$ kết quả thuận lợi.

$$P = \frac{5}{36}$$

\textbf{Nhận xét.} Số cặp có tổng bằng $k$ giảm dần khi $k$ rời xa $7$. Tổng $7$ có $6$ cặp, tổng $8$ có $5$ cặp, tổng $9$ có $4$ cặp.
\end{dapan}

\begin{dapan}[5 — Đáp án B]
\textbf{Giải.} Tung đồng xu $3$ lần cho $2^3 = 8$ kết quả. Số cách chọn vị trí cho $2$ lần mặt ngửa là $C_3^2 = 3$.

$$P = \frac{C\_3^2}{2^3} = \frac{3}{8}$$

\textbf{Chú ý.} Đây là công thức Bernoulli với $n = 3$, $k = 2$, $p = \dfrac{1}{2}$. Hệ số $C_3^2$ là phần học sinh hay quên, dẫn đến chọn A.
\end{dapan}

\begin{dapan}[6 — Đáp án A]
\textbf{Giải.} Lấy lần lượt không hoàn lại $2$ sản phẩm. Lần thứ nhất có $95$ chính phẩm trong $100$ sản phẩm.

Sau khi lấy ra một chính phẩm, còn $94$ chính phẩm trong $99$ sản phẩm.

$$P = \frac{95}{100} \times \frac{94}{99}$$

\textbf{Bẫy.} Đáp án C giữ nguyên mẫu số $100$ ở lần thứ hai, tức nhầm với trường hợp có hoàn lại. Đáp án B và D tính xác suất lấy được cả hai phế phẩm.
\end{dapan}

\begin{dapan}[7 — Đáp án C]
\textbf{Giải.} Trong $20$ số từ $1$ đến $20$ có $10$ số chẵn. Số cách rút $3$ thẻ đều là số chẵn:

$$C\_{10}^3 = \frac{10 \cdot 9 \cdot 8}{6} = 120$$

Tổng số cách rút $3$ thẻ:

$$C\_{20}^3 = \frac{20 \cdot 19 \cdot 18}{6} = 1140$$

$$P = \frac{C\_{10}^3}{C\_{20}^3} = \frac{120}{1140} = \frac{10 \cdot 9 \cdot 8}{20 \cdot 19 \cdot 18} = \frac{2}{19}$$

\textbf{Bẫy.} Đáp án A nhân thêm $2$ ở tử số. Đáp án B dùng chỉnh hợp $P_{20}^3$ ở mẫu số, tức tính đến thứ tự trong khi đề rút đồng thời.
\end{dapan}

\subsection*{Nhóm B — Biến cố đối và bài toán ít nhất}

\begin{dapan}[8 — Đáp án C]
\textbf{Giải.} Gọi $A$ là biến cố “có ít nhất một quả đen”. Biến cố đối $\overline{A}$ là “cả ba quả đều trắng”.

Hộp có $4$ quả trắng và $6$ quả đen, tổng $10$ quả.

$$P(\overline{A}) = \frac{C\_4^3}{C\_{10}^3} = \frac{4}{120}$$

$$P(A) = 1 - \frac{4}{120} = \frac{116}{120}$$

\textbf{Bẫy.} Đáp án B là $\dfrac{4}{120}$, tức tính xác suất cả ba quả đều trắng rồi dừng lại, quên lấy $1$ trừ. Đáp án D là $\dfrac{119}{120}$, ứng với bài toán “ít nhất một quả trắng”.
\end{dapan}

\begin{dapan}[9 — Đáp án A]
\textbf{Giải.} Gọi $A$ là biến cố “có ít nhất một sản phẩm đạt chuẩn”. Biến cố đối $\overline{A}$ là “cả hai sản phẩm đều không đạt chuẩn”.

Tỉ lệ không đạt chuẩn là $1 - 0.98 = 0.02$. Vì hai lần lấy độc lập:

$$P(\overline{A}) = 0.02^2$$

$$P(A) = 1 - 0.02^2$$

\textbf{Bẫy.} Đáp án D là $1 - 0.98^2$, tức lấy đối của biến cố “cả hai đều đạt chuẩn”. Biến cố đối đúng phải là “cả hai đều không đạt”.
\end{dapan}

\begin{dapan}[10 — Đáp án B]
\textbf{Giải.} Gọi $A$ là “đạt môn thứ nhất”, $B$ là “đạt môn thứ hai”. Hai biến cố độc lập nên:

$$P(A \cup B) = 1 - P(\overline{A}) \cdot P(\overline{B}) = 1 - 0.2 \times 0.3 = 1 - 0.06 = 0.94$$

\textbf{Cách nhanh.} Công thức $1 - P(\overline{A}) P(\overline{B})$ cho đáp án trong một bước, không cần tìm phần giao rồi áp công thức cộng.

\textbf{Bẫy.} Đáp án A là $0.56$, tức tích $0.8 \times 0.7$, ứng với biến cố “đạt cả hai môn”. Đáp án D là $1.5$, vượt quá $1$ nên loại được ngay mà không cần tính.
\end{dapan}

\subsection*{Nhóm C — Công thức cộng và biến cố độc lập}

\begin{dapan}[11 — Đáp án A]
\textbf{Giải.} Đề cho 只允许中一个奖项 (chỉ được trúng một giải), tức ba biến cố trúng thưởng xung khắc từng đôi. Áp dụng công thức cộng cho biến cố xung khắc:

$$P = 0.01 + 0.05 + 0.1 = 0.16$$

\textbf{Bẫy.} Đáp án B là tích ba xác suất, công thức dành cho biến cố độc lập. Đáp án C là công thức của bài toán “ít nhất một”, áp dụng khi ba biến cố độc lập và có thể cùng xảy ra.
\end{dapan}

\begin{dapan}[12 — Đáp án D]
\textbf{Giải.} Gọi $A$ là biến cố “甲 hoặc 乙 được chọn”. Biến cố đối $\overline{A}$ là “cả 甲 và 乙 đều không được chọn”, tức chọn $2$ người từ $3$ người còn lại.

$$P(\overline{A}) = \frac{C\_3^2}{C\_5^2} = \frac{3}{10}$$

$$P(A) = 1 - \frac{3}{10} = \frac{7}{10}$$

\textbf{Cách nhanh.} Cách khác là cộng trực tiếp: số cách có 甲 là $C_4^1 = 4$, số cách có 乙 là $C_4^1 = 4$, trừ đi trường hợp có cả hai là $1$:

$$4 + 4 - 1 = 7 \implies P = \frac{7}{10}$$

\textbf{Chú ý.} Đáp án C là $\dfrac{3}{5}$, tức quên trừ trường hợp có cả 甲 và 乙 khi cộng trực tiếp.
\end{dapan}

\begin{dapan}[13 — Đáp án A]
\textbf{Giải.} Từ định nghĩa xác suất có điều kiện:

$$P(A \mid B) = \frac{P(A \cap B)}{P(B)} \implies P(A \cap B) = P(A \mid B) \cdot P(B)$$

Thay số:

$$P(A \cap B) = 0.7 \times 0.4 = 0.28$$

\textbf{Bẫy.} Đáp án C là $0.7$, tức chép lại xác suất có điều kiện mà không nhân với $P(B)$. Đáp án D là $0.4$, tức chép lại $P(B)$.
\end{dapan}

\subsection*{Nhóm D — Vector và số phức}

\begin{dapan}[14 — Đáp án D]
\textbf{Giải.} Khai triển tích theo hằng đẳng thức:

$$(a+b)\cdot(a-b) = a \cdot a - b \cdot b = \lvert a \rvert^2 - \lvert b \rvert^2$$

$$\lvert a \rvert^2 = 1^2 + (-5)^2 = 26, \qquad \lvert b \rvert^2 = 2^2 + 4^2 = 20$$

$$(a+b)\cdot(a-b) = 26 - 20 = 6$$

\textbf{Cách nhanh.} Không cần tính $a+b$ và $a-b$ rồi nhân. Hiệu hai bình phương phương của hai vector cho đáp án trong hai bước.
\end{dapan}

\begin{dapan}[15 — Đáp án A]
\textbf{Giải.} Tính tích vô hướng và độ dài:

$$u \cdot v = 1 \cdot \sqrt{3} + \sqrt{3} \cdot 1 = 2\sqrt{3}$$

$$\lvert u \rvert = \sqrt{1 + 3} = 2, \qquad \lvert v \rvert = \sqrt{3 + 1} = 2$$

$$\cos\theta = \frac{2\sqrt{3}}{2 \times 2} = \frac{\sqrt{3}}{2} \implies \theta = 30^\circ$$

\textbf{Chú ý.} Hai vector đối xứng nhau qua đường phân giác nên góc giữa chúng nhỏ. Đáp án C là $60^\circ$, tức nhầm với góc của vector $u$ và trục hoành.
\end{dapan}

\begin{dapan}[16 — Đáp án C]
\textbf{Giải.} Với $z = a + bi$ ta có $\overline{z} = a - bi$, do đó:

$$z \cdot \overline{z} = a^2 + b^2 = \lvert z \rvert^2$$

Từ $z \cdot \overline{z} = 9$ suy ra $\lvert z \rvert^2 = 9$, vậy $\lvert z \rvert = 3$.

\textbf{Cách nhanh.} Nhớ đẳng thức $\lvert z \rvert^2 = z \cdot \overline{z}$. Đây là cách tính môđun không cần tách phần thực và phần ảo.
\end{dapan}

\begin{dapan}[17 — Đáp án B]
\textbf{Giải.} Hai vector cùng phương khi tọa độ tương ứng tỉ lệ. Xét từng cặp:

\begin{itemize}[leftmargin=1.4em, itemsep=2pt]
\item \textbf{A.} $(1,2)$ và $(4,4)$: tỉ số $\dfrac{4}{1} = 4 \neq 2 = \dfrac{4}{2}$. Không cùng phương.
\item \textbf{B.} $(1,2)$ và $(2,4)$: $\dfrac{2}{1} = 2 = \dfrac{4}{2}$. Cùng phương.
\item \textbf{C.} $(3,0)$ và $(0,1)$: vector thứ hai có hoành độ $0$, tỉ số không tồn tại và tích chéo $3 \cdot 1 - 0 \cdot 0 = 3 \neq 0$. Không cùng phương.
\item \textbf{D.} $(1,1)$ và $(-1,1)$: tích chéo $1 \cdot 1 - 1 \cdot (-1) = 2 \neq 0$. Không cùng phương.
\end{itemize}


\textbf{Cách nhanh.} Dùng tích chéo $a_1 b_2 - a_2 b_1 = 0$ làm điều kiện cùng phương, tránh chia cho $0$ ở cặp C.
\end{dapan}

\begin{dapan}[18 — Đáp án D]
\textbf{Giải.} Từ định nghĩa vector $\overrightarrow{AB} = B - A$ suy ra $B = A + \overrightarrow{AB}$.

$$B = (1, 2, 0) + (2, -1, 3) = (3, 1, 3)$$

\textbf{Bẫy.} Đáp án B là $(-1, 3, -3)$, tức trừ thay vì cộng. Đáp án A chép lại chính vector $\overrightarrow{AB}$.
\end{dapan}

\begin{dapan}[19 — Đáp án B]
\textbf{Giải.} Từ $z(1+i) = 2i$ suy ra:

$$z = \frac{2i}{1+i} = \frac{2i(1-i)}{(1+i)(1-i)} = \frac{2i - 2i^2}{2} = \frac{2i + 2}{2} = 1 + i$$

\textbf{Cách nhanh.} Nhân cả tử và mẫu với số phức liên hợp của mẫu. Nhớ $i^2 = -1$ nên $-2i^2 = 2$.
\end{dapan}

\subsection*{Nhóm E — Conic}

\begin{dapan}[20 — Đáp án C]
\textbf{Giải.} Mẫu số lớn hơn nằm dưới $x^2$ nên $a^2 = 25$, $b^2 = 9$ và tiêu điểm nằm trên trục hoành.

$$c^2 = a^2 - b^2 = 25 - 9 = 16 \implies c = 4$$

Tiêu điểm là $(\pm 4, 0)$.

\textbf{Bẫy.} Đáp án A là $(\pm 3, 0)$, tức lấy $\sqrt{b^2} = 3$ làm $c$. Phải tính $c^2 = a^2 - b^2$ trước khi lấy căn.
\end{dapan}

\begin{dapan}[21 — Đáp án A]
\textbf{Giải.} Đưa về dạng chính tắc $y^2 = 2px$ với $2p = 20$, suy ra $p = 10$.

Tiêu điểm của parabol $y^2 = 2px$ là $\left(\dfrac{p}{2}, 0\right) = (5, 0)$.

\textbf{Bẫy.} Đáp án C là $(0,5)$, tức đặt tiêu điểm lên trục tung trong khi phương trình có dạng $y^2 = 2px$ nên trục đối xứng là trục hoành.
\end{dapan}

\begin{dapan}[22 — Đáp án C]
\textbf{Giải.} Chia cả hai vế cho $36$ để đưa về dạng chính tắc:

$$\frac{x^2}{9} + \frac{y^2}{4} = 1$$

Suy ra $a^2 = 9$, tức $a = 3$. Độ dài trục lớn là $2a = 6$.

\textbf{Bẫy.} Đáp án A là $3$, tức trả lời độ dài nửa trục lớn thay vì trục lớn. Đề hỏi 长轴长 (độ dài trục lớn), không phải 半长轴 (nửa trục lớn).
\end{dapan}

\begin{dapan}[23 — Đáp án B]
\textbf{Giải.} Với hyperbol $\dfrac{x^2}{a^2} - \dfrac{y^2}{b^2} = 1$, hai đường tiệm cận là $y = \pm\dfrac{b}{a}x$.

Ở đây $a^2 = 16$ nên $a = 4$, $b^2 = 9$ nên $b = 3$.

$$y = \pm\frac{3}{4}x$$

\textbf{Cách nhanh.} Cho vế phải của phương trình chính tắc bằng $0$ rồi giải ra $y$, không cần nhớ công thức:

$$\frac{x^2}{16} - \frac{y^2}{9} = 0 \implies y = \pm\frac{3}{4}x$$

\textbf{Bẫy.} Đáp án A đảo ngược tỉ số thành $\dfrac{a}{b}$. Chú ý mẫu số dưới $x^2$ đi vào tử số của hệ số góc.
\end{dapan}

\subsection*{Nhóm F — Ôn tập tích lũy}

\begin{dapan}[24 — Đáp án C]
\textbf{Giải.} Với $x > 2$ thì $x - 2 > 0$. Tách số hạng và áp dụng bất đẳng thức Cauchy:

$$x + \frac{1}{x-2} = (x-2) + \frac{1}{x-2} + 2 \ge 2\sqrt{(x-2) \cdot \frac{1}{x-2}} + 2 = 2 + 2 = 4$$

Dấu bằng xảy ra khi $x - 2 = \dfrac{1}{x-2}$, tức $x - 2 = 1$, hay $x = 3$.

\textbf{Chú ý.} Không áp dụng Cauchy trực tiếp cho $x$ và $\dfrac{1}{x-2}$ vì hai số hạng không cùng biến. Phải tách bớt $2$ để hai số hạng còn lại có tích không đổi.
\end{dapan}

\begin{dapan}[25 — Đáp án D]
\textbf{Giải.} Phân tích thành nhân tử:

$$x^3 - 4x = x(x^2 - 4) = x(x-2)(x+2) = 0$$

Phương trình có ba nghiệm phân biệt $x = 0$, $x = 2$, $x = -2$. Vậy tập hợp $P$ có $3$ phần tử.

\textbf{Bẫy.} Đáp án C là $2$, tức chỉ giải $x^2 - 4 = 0$ và bỏ sót nghiệm $x = 0$. Luôn đặt nhân tử chung $x$ ra trước khi giải.
\end{dapan}

\subsection*{Nhóm F — Ôn tập tích lũy}

\begin{dapan}[1 — Đáp án D]
\textbf{Giải.} Với phân phối chuẩn $\xi \sim N(\mu, \sigma^2)$, đường cong đối xứng qua $x = \mu$.

Hai mệnh đề 乙 và 丙 cùng đúng hoặc cùng sai, vì $P(\xi > a) = 0.5$ tương đương $P(\xi \le a) = 0.5$. Đề cho chỉ có một mệnh đề sai nên 乙 và 丙 đều đúng, suy ra $a = \mu$.

Với $a = \mu$, xét mệnh đề 甲. Theo tính đối xứng:

$$P(\xi > a + 2) = P(\xi < a - 2)$$

Vì $a - 1 > a - 2$ nên $P(\xi < a - 1) > P(\xi < a - 2) = P(\xi > a + 2)$. Vậy 甲 đúng.

Xét mệnh đề 丁: khoảng $(a, a+1)$ nằm gần tâm hơn khoảng $(a+1, a+2)$, nên diện tích dưới đường cong trên khoảng thứ nhất lớn hơn.

$$P(a < \xi < a+1) > P(a+1 < \xi < a+2)$$

Mệnh đề 丁 phát biểu ngược lại nên 丁 sai.

\textbf{Nhận xét.} Bước quyết định là nhận ra 乙 và 丙 cùng giá trị chân lí. Nếu bỏ qua điều này thì không xác định được vị trí của $a$.
\end{dapan}

\begin{dapan}[2 — Đáp án C]
\textbf{Giải.} Chọn $3$ đèn để tắt trong $5$ đèn tương đương chọn $2$ đèn để sáng. Tổng số cách chọn:

$$C\_5^2 = 10$$

Đếm số cách chọn $2$ đèn sáng kề nhau trong $5$ vị trí xếp thành hàng:

$$(1,2), (2,3), (3,4), (4,5)$$

Có $4$ cặp kề nhau. Vậy số cách chọn $2$ đèn sáng không kề nhau là $10 - 4 = 6$.

$$P = \frac{6}{10} = 0.6$$

\textbf{Cách nhanh.} Với $n$ vị trí xếp hàng, số cặp kề nhau luôn là $n - 1$. Không cần liệt kê khi $n$ lớn.
\end{dapan}

\begin{dapan}[3 — Đáp án B]
\textbf{Giải.} Phân phối chuẩn $\xi \sim N(1, \sigma^2)$ đối xứng qua $\mu = 1$. Từ $P(\xi \le 0) = P(\xi \ge a)$ suy ra $0$ và $a$ đối xứng qua $1$:

$$\frac{0 + a}{2} = 1 \implies a = 2$$

Cần tìm giá trị nhỏ nhất của $f(x) = \dfrac{1}{x} + \dfrac{9}{2-x}$ trên $0 < x < 2$. Áp dụng bất đẳng thức Cauchy dạng cộng mẫu:

$$\left(\frac{1}{x} + \frac{9}{2-x}\right)\big(x + (2-x)\big) \ge (1 + 3)^2 = 16$$

Vì $x + (2-x) = 2$ nên:

$$\frac{1}{x} + \frac{9}{2-x} \ge \frac{16}{2} = 8$$

Dấu bằng xảy ra khi $\dfrac{1}{x} = \dfrac{3}{2-x}$, tức $2 - x = 3x$, hay $x = \dfrac{1}{2}$ (thỏa mãn $0 < x < 2$).

\textbf{Cách nhanh.} Công thức tổng quát: $\dfrac{m^2}{u} + \dfrac{n^2}{v} \ge \dfrac{(m+n)^2}{u+v}$ với $u, v > 0$. Ở đây $m = 1$, $n = 3$, $u + v = 2$.
\end{dapan}

\begin{dapan}[4 — Đáp án D]
\textbf{Giải.} Xếp $5$ chữ khác nhau thành một hàng có $5! = 120$ cách.

Đếm số cách xếp mà ba chữ 仁, 义, 礼 giữ nguyên thứ tự tương đối. Với mỗi cách xếp ba chữ này theo đúng thứ tự 仁 trước 义 trước 礼, hai chữ còn lại 智 và 信 nhận hai vị trí bất kì:

$$\frac{5!}{3!} = \frac{120}{6} = 20$$

$$P = \frac{20}{120} = \frac{1}{6}$$

\textbf{Cách nhanh.} Trong $3! = 6$ hoán vị của ba chữ, chỉ đúng một hoán vị giữ nguyên thứ tự 仁, 义, 礼. Vậy đáp án là $\dfrac{1}{3!} = \dfrac{1}{6}$.
\end{dapan}

\begin{dapan}[5 — Đáp án A]
\textbf{Giải.} Tính trực tiếp theo hai trường hợp cho lần rút của 甲.

Nếu 甲 rút được quả đỏ (xác suất $\dfrac{2}{5}$), còn $1$ đỏ trong $4$ quả:

$$P\_1 = \frac{2}{5} \times \frac{1}{4} = \frac{1}{10}$$

Nếu 甲 rút được quả trắng (xác suất $\dfrac{3}{5}$), còn $2$ đỏ trong $4$ quả:

$$P\_2 = \frac{3}{5} \times \frac{2}{4} = \frac{3}{10}$$

$$P = P\_1 + P\_2 = \frac{1}{10} + \frac{3}{10} = \frac{2}{5}$$

\textbf{Cách nhanh.} Không cần chia trường hợp. Rút không hoàn lại nhưng thứ tự không ảnh hưởng đến phân phối xác suất của từng vị trí, nên xác suất 乙 rút được quả đỏ vẫn là $\dfrac{2}{5}$, đúng bằng xác suất của 甲.

\textbf{Nhận xét.} Kết quả này đúng cho mọi người rút, không chỉ người thứ hai. Đây là tính chất đối xứng của phép rút không hoàn lại.
\end{dapan}

\begin{dapan}[6 — Đáp án C]
\textbf{Giải.} Không gian mẫu là $\{1,2,3,4,5,6\}$.

$$A = \{1, 4\}, \qquad B = \{1, 3, 5\}$$

Vì $A \cap B = \{1\} \neq \varnothing$ nên $A$ và $B$ không xung khắc, do đó đáp án A sai.

Vì $A \cup B = \{1,3,4,5\} \neq \Omega$ nên $A$ và $B$ không đối nhau, đáp án B sai.

$$P(A + B) = P(A \cup B) = \frac{4}{6} = \frac{2}{3}$$

\textbf{Chú ý.} Kí hiệu $A + B$ trong tài liệu Trung Quốc nghĩa là biến cố hợp $A \cup B$, không phải phép cộng số học.
\end{dapan}

\begin{dapan}[7 — Đáp án D]
\textbf{Giải.} Mỗi lần rút có $2$ kết quả có thể (đỏ hoặc đen) và có hoàn lại, nên ba lần rút độc lập nhau.

Số điểm mẫu là $2^3 = 8$.

\textbf{Cách nhanh.} Với $k$ lần rút có hoàn lại từ $m$ loại kết quả, số điểm mẫu là $m^k$. Ở đây $2^3 = 8$.

\textbf{Bẫy.} Đáp án B là $6$, ứng với việc đếm số lượng từng màu (từ $0$ đến $3$ quả đỏ) thay vì đếm dãy kết quả có thứ tự.
\end{dapan}

\begin{dapan}[8 — Đáp án B]
\textbf{Giải.} Gọi $A$ là biến cố “có ít nhất một lần trúng thưởng”. Biến cố đối $\overline{A}$ là “cả ba lần đều không trúng”.

Xác suất không trúng mỗi lần là $1 - 0.1 = 0.9$. Ba lần độc lập nên:

$$P(\overline{A}) = 0.9^3$$

$$P(A) = 1 - 0.9^3$$

\textbf{Bẫy.} Đáp án D là $1 - 0.1^3$, tức lấy đối của biến cố “cả ba lần đều trúng”. Biến cố đối đúng là “không lần nào trúng”.
\end{dapan}

\begin{dapan}[9 — Đáp án A]
\textbf{Giải.} Lần thứ nhất lấy từ $10$ quả, trong đó có $4$ quả xanh:

$$P(\text{lần 1 xanh}) = \frac{4}{10}$$

Sau khi lấy ra một quả xanh và không hoàn lại, còn $9$ quả trong đó có $6$ quả trắng:

$$P(\text{lần 2 trắng} \mid \text{lần 1 xanh}) = \frac{6}{9}$$

$$P = \frac{4}{10} \times \frac{6}{9}$$

\textbf{Bẫy.} Đáp án B giữ nguyên mẫu số $10$ ở lần hai, tức nhầm với trường hợp có hoàn lại. Đáp án C cộng hai xác suất thay vì nhân.
\end{dapan}

\begin{dapan}[10 — Đáp án B]
\textbf{Giải.} Vì $A$ và $B$ xung khắc nên $P(A \cap B) = 0$. Áp dụng công thức cộng:

$$P(A \cup B) = 0.6 + 0.3 = 0.9$$

\textbf{Bẫy.} Đáp án D là $0.18$, tức tích hai xác suất, công thức dành cho biến cố độc lập. Đáp án A là $0.3$, chỉ là $P(B)$.
\end{dapan}

\begin{dapan}[11 — Đáp án A]
\textbf{Giải.} Gọi $F$ là “thích bóng đá”, $B$ là “thích bóng rổ”. Đề cho:

$$P(F) = 0.6, \qquad P(B) = 0.3, \qquad P(F \cap B) = 0.1$$

Điều kiện đã biết là em đó thích bóng rổ, nên mẫu số là $P(B)$:

$$P(F \mid B) = \frac{P(F \cap B)}{P(B)} = \frac{0.1}{0.3} = \frac{1}{3}$$

\textbf{Bẫy.} Đáp án C là $\dfrac{1}{6}$, tức chia cho $P(F) = 0.6$. Mẫu số phải là xác suất của nhóm đứng sau chữ 如果 (nếu).
\end{dapan}

\begin{dapan}[12 — Đáp án C]
\textbf{Giải.} Hai biến cố $\{X > a\}$ và $\{X \le a\}$ là hai biến cố đối, nên tổng xác suất bằng $1$:

$$P(X \le a) = 1 - P(X > a) = 1 - 0.3 = 0.7$$

\textbf{Cách nhanh.} Đây là bài toán biến cố đối thuần túy, không cần biết phân phối của $X$ là gì.

\textbf{Bẫy.} Đáp án D là 无法确定 (không xác định được), nhưng thực ra xác định được vì quan hệ đối luôn đúng với mọi biến ngẫu nhiên.
\end{dapan}

\begin{dapan}[13 — Đáp án C]
\textbf{Giải.} Không gian mẫu là $\{1,2,3,4,5,6\}$.

$$A = \{2, 4, 6\}, \qquad B = \{5, 6\}$$

$$A \cap B = \{6\}, \qquad A \cup B = \{2, 4, 5, 6\}$$

Áp dụng công thức cộng tổng quát:

$$P(A \cup B) = \frac{3}{6} + \frac{2}{6} - \frac{1}{6} = \frac{4}{6} = \frac{2}{3}$$

\textbf{Bẫy.} Đáp án D là $\dfrac{5}{6}$, tức quên trừ phần giao. Hai biến cố này có phần tử chung là $6$ nên không xung khắc.
\end{dapan}

\begin{dapan}[14 — Đáp án A]
\textbf{Giải.} Gọi $A$ là biến cố “có ít nhất $2$ nam”. Biến cố đối $\overline{A}$ là “có ít nhất $4$ nữ”, tức chỉ có $0$ hoặc $1$ nam.

Đếm số cách thuận lợi cho $\overline{A}$: chọn $5$ nữ, hoặc chọn $1$ nam và $4$ nữ.

$$P(\overline{A}) = \frac{C\_{15}^5 + C\_{10}^1 C\_{15}^4}{C\_{25}^5}$$

$$P(A) = 1 - \frac{C\_{15}^5 + C\_{10}^1 C\_{15}^4}{C\_{25}^5}$$

\textbf{Bẫy.} Đáp án B và D chỉ tính trường hợp có đúng $2$ nam, bỏ sót trường hợp $3$, $4$, $5$ nam. Đáp án C lấy đối của biến cố “có ít nhất $2$ nam” nhưng phần bù lại đếm các trường hợp có $2$ nam trở lên.
\end{dapan}

\begin{dapan}[15 — Đáp án B]
\textbf{Giải.} Mức ý nghĩa $\alpha$ là xác suất bác bỏ giả thuyết $H_0$ trong khi $H_0$ đúng. Đây chính là định nghĩa của sai lầm loại một.

Với $\alpha = 0.05$:

$$P(\text{sai lầm loại một}) = 0.05 = 5\%$$

\textbf{Bẫy.} Đáp án A là $95\%$, tức độ tin cậy $1 - \alpha$ chứ không phải mức ý nghĩa. Hai con số này bù nhau nên rất dễ chọn nhầm.
\end{dapan}

\clearpage
\sectiondautrang{Tổng kết lỗi thường gặp}
\addcontentsline{toc}{section}{Tổng kết lỗi thường gặp}

Bảng chẩn đoán dùng khi chữa bài và khi soạn đề kiểm tra. Cột ``Câu liên quan'' trỏ tới đề bài tập về nhà.

\bang{
\begin{tabular}{@{}p{5.6cm}p{4.2cm}p{5.2cm}@{}}
\toprule
\textbf{Lỗi} & \textbf{Câu liên quan} & \textbf{Cách cảnh báo học sinh} \\
\midrule
Dùng công thức xung khắc cho biến cố độc lập, hoặc ngược lại & Câu 11, 13 phần bắt buộc; Câu 10 phần tự chọn & Viết to 互斥 và 独立 lên bảng: xung khắc thì phần giao bằng $0$, độc lập thì phần giao bằng tích \\
Quên hệ số $C_n^k$ ở công thức Bernoulli & Câu 5 phần bắt buộc & $k$ lần thành công rơi vào vị trí nào cũng được, nên luôn có hệ số đếm vị trí \\
Lấy $1$ trừ sai biến cố đối & Câu 8, 9, 10 phần bắt buộc; Câu 8, 14 phần tự chọn & Yêu cầu viết biến cố đối bằng lời trước khi lấy $1$ trừ \\
Giữ nguyên mẫu số ở lần lấy thứ hai khi không hoàn lại & Câu 6 phần bắt buộc; Câu 9 phần tự chọn & Đề không ghi 有放回 thì mẫu số phải giảm sau mỗi lần lấy \\
Chia nhầm mẫu số ở bài xác suất có điều kiện & Câu 11 phần tự chọn & Khoanh tròn cụm sau chữ 如果: mẫu số là nhóm đó, không phải cả không gian mẫu \\
Dùng chỉnh hợp thay vì tổ hợp khi đếm đồng thời & Câu 7 phần bắt buộc & Đổi chỗ hai vật đã chọn có ra kết quả khác không, không thì dùng $C$ \\
Nhầm $C_n^2$ với tích $n(n-1)$ ở mẫu số & Câu 3, 4, 5 phần bắt buộc; Câu 2, 5 phần tự chọn & Rút hai vật cùng lúc: mẫu số luôn là $\dfrac{n(n-1)}{2}$ \\
Nhầm biến cố đối với biến cố bù một phần & Câu 9 phần bắt buộc & ít nhất một đạt đối với cả hai đều không đạt, không phải cả hai đều đạt \\
Đọc nhầm 至少 và 恰好 & Câu 8 phần bắt buộc; Câu 5 phần tự chọn & Ít nhất thì đi qua biến cố đối, đúng $k$ lần thì đếm trực tiếp \\
Chọn nhầm trục của tiêu điểm conic & Câu 20 phần bắt buộc & Mẫu lớn hơn nằm dưới biến nào thì tiêu điểm nằm trên trục của biến kia \\
Lấy $\sqrt{b^2}$ làm $c$ của elip & Câu 20 phần bắt buộc & Phải tính $c^2 = a^2 - b^2$ trước khi lấy căn \\
Nhầm trục lớn với nửa trục lớn & Câu 22 phần bắt buộc & Đề hỏi 长轴长 thì nhân đôi $a$, đừng dừng ở $a$ \\
Bỏ sót nghiệm khi phân tích đa thức & Câu 25 phần bắt buộc & Đặt nhân tử chung $x$ ra trước, rồi mới giải phương trình bậc hai còn lại \\
Cộng trực tiếp mà quên trừ phần giao & Câu 12 phần bắt buộc; Câu 13 phần tự chọn & Hai biến cố có phần tử chung thì phải trừ, kể cả khi chỉ chung một phần tử \\
\bottomrule
\end{tabular}}

% Không dùng \vfill: nó đẩy chữ xuống sát đáy trang,
% nơi banner letterhead chiếm chỗ (281-297mm) và sẽ che mất chữ.
\vspace{0.8cm}
\begin{center}
{\footnotesize\color{cae_blue} Tài liệu nội bộ, CAE SHANGHAI}
\end{center}

\end{document}
